\chapter{Results and Discussion}

\section{Output of the Working System}

The application is usable end to end against the composed backend. The screens below
show the driver's path from discovery to a closed history entry; every value on them
comes from the backend, not from a fixture.

\reportimage[0.55]{\ImgScreenDashboard}{Dashboard after sign-in: the active session
and the vehicles held by the signed-in account.}{fig:screen-dashboard}

\reportimage[0.55]{\ImgScreenDiscover}{Discovery: the gate advertising
\code{DWARA-FAC-0001} appears in the scan, and its detail is fetched from the backend
by facility code.}{fig:screen-discover}

\reportimage[0.55]{\ImgScreenActive}{Active session: the entry gate button, the live
status and the running fee.}{fig:screen-active}

\reportimage[0.55]{\ImgScreenPayment}{Payment: the frozen amount after a successful
mock payment, with the exit gate now available.}{fig:screen-payment}

\reportimage[0.55]{\ImgScreenActivity}{Activity: the finished session listed from the
paged history endpoint.}{fig:screen-activity}

\subsection{End-to-end walkthrough}
\label{sec:walkthrough}

The walkthrough below was executed against the development stack with a seeded demo
facility and one registered gate. Each step lists the observable result rather than
the implementation behind it.

\begin{table}[htbp]
  \centering
  \small
  \begin{tabularx}{\textwidth}{@{}clX@{}}
    \toprule
    \textbf{Step} & \textbf{Action} & \textbf{Observed result} \\
    \midrule
    1 & Register, then sign in & Account created; token stored; dashboard shown \\
    2 & Add a vehicle & Plate and class accepted; vehicle appears on the dashboard \\
    3 & Scan for gates & \code{DWARA-FAC-0001} listed; facility detail shows the
      per-second rates \\
    4 & Start a session & \code{CREATED}; the identifier is shown and the entry
      button becomes available \\
    5 & Open the entry gate & Backend returns \code{open}; the servo sweeps; status
      becomes \code{PARKED} with an entry timestamp \\
    6 & Wait and watch the fee & Amount increases once per second at the facility
      rate \\
    7 & Pay & Receipt with a \code{MOCK-} reference; amount frozen; session
      \code{PAYMENT\_PENDING} \\
    8 & Open the exit gate & Backend verifies the payment; gate opens; session
      \code{EXITED} \\
    9 & Open the activity screen & The finished session appears with its frozen
      total \\
    10 & Queue a remote \code{OPEN} (operator) & Gate executes on its next poll and
      reports \code{sessionId: "remote"} \\
    \bottomrule
  \end{tabularx}
  \caption{End-to-end walkthrough.}
  \label{tab:walkthrough}
\end{table}

Two behaviours in the walkthrough deserve emphasis because they demonstrate that the
architecture is doing real work. In step 5 the application cannot open the gate: it
only delivers the request, and the state change is visible because the backend
recorded it. In step 8 the exit is granted by payment state, not by the phone, so the
same session cannot be used to re-enter the facility.

\section{Analysis Against Objectives}

\begin{table}[htbp]
  \centering
  \small
  \begin{tabularx}{\textwidth}{@{}lXl@{}}
    \toprule
    \textbf{Objective} & \textbf{Evidence} & \textbf{Outcome} \\
    \midrule
    Frictionless discovery & The gate advertises \code{DWARA-<facilityCode>}; the
      application lists it from a scan and resolves details through
      \code{GET /facilities/\{code\}}. & Met \\
    Server-authoritative control & Every entry and exit is a
      \code{POST /gate/validate} decision; the firmware has no rule that can open a
      gate on its own, and fails closed on transport errors. & Met \\
    Auditable per-second billing & Amounts are computed from
      \code{entryTimestamp} and the facility rate, frozen at payment, and returned
      with a receipt; the client never supplies a time or an amount. & Met \\
    Safe lifecycle & One active session per user and per vehicle, guarded by a row
      lock; stale sessions reaped; every illegal transition has a reason code and a
      test. & Met \\
    Three-surface deliverable & Backend, Expo application and ESP32 firmware in one
      repository, with a shared configuration model and one Compose stack. & Met \\
    Reproducible verification & One command builds, boots, tests and publishes
      reports; 118 tests pass under a 100\% service-line coverage gate. & Met \\
    \bottomrule
  \end{tabularx}
  \caption{Objectives against evidence.}
  \label{tab:objectives}
\end{table}

The only objectives that are met at the software boundary rather than on hardware are
the firmware behaviours. The sketch compiles for the ESP32 target and the serial
bypass allows the full backend flow to be exercised without Wi-Fi, but the servo
sweep, the stall watchdog and the BLE round trip on a physical gate are documented
bench procedures, not results measured in this environment. That is a limitation of
the evaluation, not of the design, and it is the first item in the future-work list.

\subsection*{Measured observations}

Table~\ref{tab:verification} in the previous chapter records the verification
results. Three numbers are worth reading as outcomes rather than as counts:

\begin{itemize}
  \item \textbf{118 automated tests, 0 failures, 100\% service-line coverage.} The
        coverage gate is enforced in \code{verify}, so the figure cannot drift
        silently.
  \item \textbf{Sub-second backend responses in the local stack.} The authentication
        contract run recorded responses between roughly 6\,ms and 155\,ms, with the
        password-hashing calls in registration and login at the slow end of that
        range.
  \item \textbf{Deterministic gate timing.} The open hold is a fixed 8\,s and the
        settle window 800\,ms, so the application's 15\,s round-trip budget leaves
        ample margin even when the backend is slow.
\end{itemize}

\subsection*{Discussion}

The most useful result is negative: the system works without any component that
identifies a vehicle. A phone account, a session identifier and a server-side state
machine were sufficient to authorise entry and exit, and the only hardware at the
site is a microcontroller and a hobby servo. That is a meaningful cost reduction for
a small facility, where a camera or an RFID reader per lane would dominate the
budget.

The second result is that centralising the decision simplified the edge. The firmware
contains no pricing, no payment state and no notion of a user; it contains a timeout,
a verdict and a servo. When the BLE client was fixed during development, the firmware
did not change; when the grace window moved from payment time to gate time, the
firmware did not change either. The cost of that simplification is a hard dependency
on the backend at the moment of entry and exit, and a deliberate one: the gate fails
closed, and the operator has a remote-open path for the post-payment case.

The third result concerns testability. Because the payment provider, the gate client
and the actuator are interfaces, the rule-heavy code is testable without a database
or a servo. The classes with the most tests are the two that decide money and access,
which is where defects would be most expensive.

The design also has costs that should be stated plainly. A fixed-timer gate cannot
detect tailgating; plaintext BLE means the session identifier is visible to anything
in range; a single gate per facility means one queue serves both directions; and the
mock payment provider means the settlement path has never seen a real failure. Each
of these is a deliberate version-one boundary, and each has a concrete replacement in
the next section.
